% =============================================================================
% table_components.tex — CUMCM 表格系统（v4.5 模板 v2 §13）
% 短表/自适应表/长表/宽表/结果表/符号表/对比表；默认 [htbp]；三线表（booktabs）。
% 禁止默认 \resizebox{\textwidth}{!}（宽表优先级：重写列名 → 自适应列宽 → 局部 small → landscape）。
% =============================================================================

% ---------- \ShortTable{标题}{列定义}{表头}{行}：短文本表，table + tabularx，默认 [htbp] ----------
\newcommand{\ShortTable}[4]{%
  \begin{table}[htbp]
    \centering
    \caption{#1}
    \vspace{0.3em}
    \begin{tabularx}{\textwidth}{#2}
      \toprule
      #3 \\
      \midrule
      #4 \\
      \bottomrule
    \end{tabularx}
  \end{table}%
}

% ---------- \FlexibleTable{标题}{列定义}{表头}{行}：X/p 列自动换行的自适应宽表 ----------
% 列定义示例：{@{}l X r@{}}（X 列自动填充剩余宽度并换行）
\newcommand{\FlexibleTable}[4]{%
  \begin{table}[htbp]
    \centering
    \caption{#1}
    \vspace{0.3em}
    \begin{tabularx}{\textwidth}{#2}
      \toprule
      #3 \\
      \midrule
      #4 \\
      \bottomrule
    \end{tabularx}
  \end{table}%
}

% ---------- \LongTable{标题}{列定义}{表头}{行}：longtable 跨页（符号表/文件清单/长结果表） ----------
% TPL-L02：四参数命令（禁止再以 \begin{LongTable} 环境方式使用）；
% TPL-L03：longtable 不支持 tabularx 的 X 列——列定义必须用 l/p{}/等固定宽度；
% TPL-L04：body 行由生成器自行以 \\ 结束，组件不再统一追加 \\（空 body 合法）。
\newcommand{\LongTable}[4]{%
  \begingroup
  \small
  \setlength{\tabcolsep}{4pt}
  \renewcommand{\arraystretch}{1.15}
  \setlength{\LTcapwidth}{\textwidth}
  \begin{longtable}{#2}
    \caption{#1} \\
    \toprule
    #3 \\
    \midrule
    \endfirsthead
    \toprule
    #3 \\
    \midrule
    \endhead
    #4
    \bottomrule
  \end{longtable}
  \endgroup
}

% ---------- \ResultTable{标题}{列定义}{表头}{行}：一问最终结果表；答案格用 \resultcell ----------
\newcommand{\ResultTable}[4]{%
  \begin{table}[htbp]
    \centering
    \caption{#1}
    \vspace{0.3em}
    \begin{tabularx}{\textwidth}{#2}
      \toprule
      #3 \\
      \midrule
      #4 \\
      \bottomrule
    \end{tabularx}
  \end{table}%
}
% \resultcell{...}：仅强调真正答案（禁止整行粗体）
\newcommand{\resultcell}[1]{{\heiti\bfseries #1}}

% ---------- \SymbolTable{标题}{行}：符号/含义/单位三列，语义固定，longtable 可跨页 ----------
% 行内容：$N$ & 样本数量 & 个 \\
% TPL-L03：longtable 禁用 X 列（X 仅 tabularx 有效）——固定 p{} 宽度；
% TPL-L04：行自行以 \\ 结束，空 body 合法（不暴力追加 \\）。
\newcommand{\SymbolTable}[2]{%
  \begingroup
  \small
  \setlength{\tabcolsep}{4pt}
  \renewcommand{\arraystretch}{1.15}
  \setlength{\LTcapwidth}{\textwidth}
  \begin{longtable}{@{}l p{8.6cm} p{2.4cm}@{}}
    \caption{#1} \\
    \toprule
    \textbf{符号} & \textbf{含义} & \textbf{单位} \\
    \midrule
    \endfirsthead
    \toprule
    \textbf{符号} & \textbf{含义} & \textbf{单位} \\
    \midrule
    \endhead
    #2
    \bottomrule
  \end{longtable}
  \endgroup
}

% ---------- \ComparisonTable{标题}{列定义}{表头}{行}：模型/方案对比表（含 BASELINE 标识） ----------
\newcommand{\ComparisonTable}[4]{%
  \begin{table}[htbp]
    \centering
    \caption{#1}
    \vspace{0.3em}
    \begin{tabularx}{\textwidth}{#2}
      \toprule
      #3 \\
      \midrule
      #4 \\
      \bottomrule
    \end{tabularx}
  \end{table}%
}

% ---------- \WideTable{标题}{列定义}{表头}{行}：超宽表——自适应列宽 + 局部 small 兜底 ----------
% 禁止默认 resizebox；列定义用 X/p 列优先，实在超宽再用 \small 或 landscape（\begin{landscape}）。
% TPL-L01：本文件是全模板唯一 \tablesmall 定义处（main.tex 不得重复定义）。
\newcommand{\tablesmall}{\small}
\newcommand{\WideTable}[4]{%
  \begin{table}[htbp]
    \centering
    \caption{#1}
    \vspace{0.3em}
    {\tablesmall
    \begin{tabularx}{\textwidth}{#2}
      \toprule
      #3 \\
      \midrule
      #4 \\
      \bottomrule
    \end{tabularx}}
  \end{table}%
}
